\documentclass[11pt]{article}
\usepackage[margin=1in]{geometry}
\usepackage{amsmath,amssymb,amsthm}
\usepackage[T1]{fontenc}
\usepackage{lmodern,microtype}
\usepackage[hidelinks]{hyperref}
\newtheorem{theorem}{Theorem}
\newtheorem{lemma}{Lemma}
\newcommand{\HC}{\operatorname{HC}}
\newcommand{\FHC}{\operatorname{FHC}}
\newcommand{\C}{\mathbb C}
\newcommand{\norm}[2]{\lVert #1\rVert_{#2}}
\title{Polynomial multipliers of differentiation-hypercyclic functions}
\author{Alec Kriebel\\
\small Independent researcher\\
\small ORCID: \href{https://orcid.org/0009-0001-9320-500X}{0009-0001-9320-500X}}
\date{17 October 2026}
\begin{document}
\maketitle
\begin{abstract}
We give a self-contained elementary negative answer to the polynomial-multiplier
question posed by Costakis in 2008. There is a dense family of entire functions
hypercyclic for differentiation such that every nonconstant complex polynomial
multiple fails to be hypercyclic; in fact its derivative values at zero escape
in modulus. The proof uses prescribed annular bounds on derivative-normalized
Taylor coefficients and relative Baire category. This is an attributed
historical-answer note, without a claim of a new first solution: Nikula's older
random-function theorem already implies a stronger frequently hypercyclic
version, as explained below.
\end{abstract}

Let $H(\C)$ have the topology of uniform convergence on compact sets, write
$Df=f'$, and put $\norm{f}{R}=\max_{|z|\le R}|f(z)|$.
A function is in $\HC(D)$ if $\{D^nf:n\ge0\}$ is dense in $H(\C)$.
Costakis's Question~5 \cite[p.~330]{Costakis} asks whether $pf$ must be in
$\HC(D)$ whenever $f\in\HC(D)$ and $p$ is a nonconstant polynomial.
The operation is the pointwise product $p(z)f(z)$.

\begin{theorem}\label{thm:main}
There is a dense subset $B\subset H(\C)$, contained in $\HC(D)$, such that
for every $f\in B$ and every nonconstant $p\in\C[z]$,
\[
                 |(pf)^{(n)}(0)|\longrightarrow\infty.
\]
Consequently the polynomial multipliers preserving every member of $\HC(D)$
are exactly the nonzero constants.
\end{theorem}

The same $f$ works for all $p$; the onset of the limit may depend on $p$.
We first prove the theorem without using random hypercyclicity.
Set
\[
 \ell_n=(n+1)^{-1/2},\qquad u_n=\log(n+e),\qquad
 X_M=\{f\in H(\C):\ell_n\le |f^{(n)}(0)|\le u_n\ (n\ge M)\}.
\]
Each $X_M$ is nonempty and closed in the complete metrizable space $H(\C)$.
Indeed the baseline series $\sum_{n\ge0}\ell_nz^n/n!$ is entire, and each
derivative evaluation is continuous. The bounds $u_n\le n+1$ give normally
convergent series and uniformly small factorial tails on each fixed disk.

\begin{lemma}\label{lem:density}
For every integer $M\ge0$, $\HC(D)\cap X_M$ is dense in $X_M$.
\end{lemma}
\begin{proof}
Use the countable target polynomials with coefficients in $\mathbb Q+i\mathbb Q$.
For such a target $q$, and integers $R,s\ge1$, the relative hitting set
\[
 U(q,R,s)=\bigcup_{N\ge0}
       \{f\in X_M:\norm{D^Nf-q}{R}<1/s\}
\]
is open in $X_M$. We show it is relatively dense.

Fix $g\in X_M$ and a neighborhood $\norm{f-g}{r}<\varepsilon$.
Choose $L\ge M$ with
\[
                 2\sum_{n>L}u_n\frac{r^n}{n!}<\varepsilon.
\]
Write $q(z)=\sum_{k=0}^d\beta_kz^k/k!$ (for $q=0$, take $d=0$).
For a large $N>L$, retain the jets of $g$ through $L$, use the baseline
$\ell_n$ afterward, and replace jets $N,\ldots,N+d$ by $\gamma_k$.
Choose $\gamma_k=\beta_k$ if $|\beta_k|\ge\ell_{N+k}$ and otherwise
move $\beta_k$ radially to modulus $\ell_{N+k}$, choosing the positive
boundary point when $\beta_k=0$. Then
\[
 |\gamma_k-\beta_k|\le\ell_{N+k},\qquad
 \ell_{N+k}\le|\gamma_k|\le\max(1,|\beta_k|).
\]
Since $u_n\to\infty$, these finitely many jets satisfy their upper bounds
for sufficiently large $N$. The resulting entire function $F_N$ lies in
$X_M$ and is within $\varepsilon$ of $g$ on the $r$-disk: both tails
after $L$ have jets bounded by $u_n$.

Under $D^N$, all retained low jets vanish. On the target block the jet errors
are at most $\ell_{N+k}$, and after that block the jets are the baseline.
Thus
\[
                   \norm{D^NF_N-q}{R}\le \ell_Ne^R\longrightarrow0.
\]
Every relative neighborhood therefore meets $U(q,R,s)$.
Baire's theorem, applied inside the complete space $X_M$, gives a dense
intersection of these hitting sets. Membership means density of the derivative
orbit, since the target polynomials are dense and integer disks and
reciprocal-integer errors suffice. No invariance of $X_M$ under $D$ is needed.
\end{proof}

\begin{proof}[Proof of Theorem~\ref{thm:main}]
Let $f\in X_M$, $a_n=f^{(n)}(0)$, and
$p(z)=\sum_{j=0}^dc_jz^j$ with $d\ge1$ and $c_d\ne0$.
For $n\ge M+d$, Leibniz gives
\[
 (pf)^{(n)}(0)=\sum_{j=0}^dc_j(n)_j a_{n-j},
 \qquad (n)_j=\frac{n!}{(n-j)!}.
\]
The top-degree term has modulus at least
\[
               |c_d|\frac{(n)_d}{\sqrt{n-d+1}}
                    \asymp_p n^{d-1/2},
\]
whereas the sum of the moduli of all lower terms is at most
\[
             \left(\sum_{j<d}|c_j|\right)n^{d-1}\log(n+e)
                         =o(n^{d-1/2}).
\]
The reverse triangle inequality proves escape, including $d=1$ and arbitrary
complex phases. This deterministic argument holds for every $p$ after $f$
has been fixed; no uncountable intersection is used.

Evaluation at zero is continuous and onto $\C$. An escaping scalar sequence
has a tail outside the unit disk; an open subdisk avoiding its finitely many
earlier values remains. It is not dense, so $pf\notin\HC(D)$.

To see ambient density, approximate any entire target by a polynomial $P$.
For large $M>\deg P$, the function
$P+\sum_{n\ge M}\ell_nz^n/n!$ belongs to $X_M$ and is arbitrarily close to $P$.
Lemma~\ref{lem:density} then gives nearby hypercyclic members of $X_M$.
Thus $B=\bigcup_M(X_M\cap\HC(D))$ is dense and has the required property.
Nonzero constants preserve a dense derivative orbit by scaling; the zero
polynomial and every nonconstant polynomial fail to preserve the full class.
\end{proof}

Each $X_M$ has empty ambient interior: it is closed and contains no
polynomial, while polynomials are dense. Hence the displayed $B$ is meagre.
This is compatible with Costakis's Theorem~6 \cite[p.~330]{Costakis}, which
gives a residual common family on which every nonzero polynomial multiple
is hypercyclic.

\paragraph{Earlier implication and priority.}
The existence conclusion of Theorem~\ref{thm:main} already follows from
Nikula's random-function theorem \cite[arXiv v1, Thm.~1, pp.~2--3]{Nikula},
even with frequent hypercyclicity. Here frequent hypercyclicity means
positive lower density of visits to every nonempty open set.
For IID standard circular complex Gaussian jets $X_n$, that theorem gives
an almost surely entire, frequently hypercyclic series
$f(z)=\sum_{n\ge0}X_nz^n/n!$. Its full-support and logarithmic-tail
hypotheses hold for the density $\pi^{-1}e^{-|z|^2}$.

Indeed, $\Pr(|X|<t)\le t^2$ and $\Pr(|X|>t)=e^{-t^2}$.
For $\alpha=3/4$, Borel--Cantelli gives, on one probability-one event,
\[
       (n+1)^{-\alpha}\le |X_n|\le\log(n+e)\quad\hbox{eventually}.
\]
Intersect with the frequent-hypercyclicity event. For any polynomial of
degree $d\ge1$, the same Leibniz estimate has leading size
$n^{d-\alpha}$ and lower-term bound $O_p(n^{d-1}\log n)$; their ratio
tends to zero. The event was chosen independently of $p$, so it supplies
all-polynomial escape simultaneously. Polynomial additions preserve both
frequent hypercyclicity and all sufficiently high jets, giving a dense
family from a single realization.

Even the exact annular existence in Lemma~\ref{lem:density} is a consequence
of that older theorem and elementary clipping. Radially project jets
$n\ge M$ into $[\ell_n,u_n]$ in modulus. Upper violations occur only finitely
often almost surely; lower corrections are at most $\ell_n$.
The correction jets $c_n$ therefore tend to zero, and their entire series
$h$ satisfies
\[
              \norm{D^nh}{R}\le e^R\sup_{k\ge n}|c_k|\longrightarrow0.
\]
Such a perturbation preserves frequent hypercyclicity by retaining visits
to smaller target neighborhoods after a finite prefix.
The clipped independent coordinates have full relative support on their
annuli (at $n=0$, Gaussian phase gives full support on the unit circle), and the common factorial-tail bound makes finite-coordinate cylinders
a neighborhood basis in $X_M$. Thus frequently hypercyclic functions are
dense in each $X_M$. Nikula is applied before clipping; the clipped laws
are not IID with full-plane support.

A still earlier general Gaussian-eigenfield route is supplied by
Bayart--Matheron \cite[arXiv v1, Prop.~3.3]{BayartMatheron}: on the unit
circle with Haar measure use $E(\lambda)(z)=e^{\lambda z}$.
Its integral operator sends Fourier basis vectors to $z^n/n!$ or zero;
the Gaussian series converges normally almost surely, has full support,
and the Fourier-decay hypothesis follows from the Riemann--Lebesgue lemma.
The resulting mixing Gaussian law gives almost-sure frequent hypercyclicity.
The supporting audit records the specialization and its limits.

We do not claim that these sources explicitly state the polynomial-multiplier
corollary, nor determine its first explicit recording or earliest underpinning.
This note provides a self-contained elementary proof of the historical
negative answer with the prior implication credited; it does not claim a
new first solution, new existence theorem, or a problem still open in 2026.
The repository's original finer fixed-polynomial null-orbit perturbation
argument is retained in the supplement, without a novelty claim.

\clearpage
\paragraph{Verification, AI use, and review status.}
The accompanying archive contains this TeX source, immutable submitted and
current mathematical editions, a reconstruction from the original question
alone, independent analytic and direct-block adversarial reviews, exact
standard-library coefficient checks and receipts, and bounded primary-source
priority audits. Finite diagnostics support coefficient identities; the
infinite-dimensional conclusions rest on the proofs.
AI tools were used extensively in research, proof development, literature
review, computation, internal adversarial verification, and writing.
This preprint is unrefereed. No external human peer review, journal acceptance,
or formal proof-assistant certification is claimed. Source bodies are not
redistributed, and historical unavailable checks are not presented as fresh
reproductions.

\begin{thebibliography}{9}
\bibitem{Costakis}
G.~Costakis, \emph{Which maps preserve universal functions?},
in Mini-Workshop: Complex Approximation and Universality,
Oberwolfach Reports \textbf{5} (2008), 328--331
(whole report 297--346), Question~5 on 330.
\href{https://doi.org/10.4171/OWR/2008/06}{doi:10.4171/OWR/2008/06}.

\bibitem{Nikula}
M.~Nikula, \emph{Frequent hypercyclicity of random entire functions for the
differentiation operator}, Complex Analysis and Operator Theory \textbf{8}
(2014), 1455--1474; online 8 October 2013.
\href{https://doi.org/10.1007/s11785-013-0328-0}{doi:10.1007/s11785-013-0328-0}.
Inspected theorem: \href{https://arxiv.org/abs/1209.6209v1}{arXiv:1209.6209v1},
repository submission 27 September 2012.
The served PDF's later internal date is separately recorded in the audit.

\bibitem{BayartMatheron}
F.~Bayart and \'E.~Matheron, \emph{Mixing operators and small subsets of the
circle}, J.~Reine Angew.~Math. \textbf{715} (2016), 75--123.
\href{https://doi.org/10.1515/crelle-2014-0002}{doi:10.1515/crelle-2014-0002}.
Inspected general theorem:
\href{https://arxiv.org/abs/1112.1289v1}{arXiv:1112.1289v1},
6 December 2011, Proposition~3.3; no backdating of a later explicit remark.
\end{thebibliography}
\end{document}
